\documentclass[aspectratio=169,10pt]{beamer}

% Shared Beamer setup for Fall 2026 course slides.
\mode<presentation>{
  \usetheme{Madrid}
  \usecolortheme{default}
}

\definecolor{CalPolyGreen}{HTML}{154734}
\definecolor{CalPolyGold}{HTML}{BD8B13}
\definecolor{DeckInk}{HTML}{1F2933}
\definecolor{DeckMuted}{HTML}{5F6B7A}

\setbeamercolor{structure}{fg=CalPolyGreen}
\setbeamercolor{palette primary}{bg=CalPolyGreen,fg=white}
\setbeamercolor{palette secondary}{bg=DeckInk,fg=white}
\setbeamercolor{palette tertiary}{bg=CalPolyGold,fg=DeckInk}
\setbeamercolor{frametitle}{bg=CalPolyGreen,fg=white}
\setbeamercolor{title}{fg=CalPolyGreen}
\setbeamercolor{normal text}{fg=DeckInk,bg=white}
\setbeamercolor{block title}{bg=CalPolyGreen,fg=white}
\setbeamercolor{block body}{bg=black!3,fg=DeckInk}

\setbeamertemplate{navigation symbols}{}
\setbeamertemplate{footline}[frame number]
\setbeamertemplate{itemize item}{\color{CalPolyGold}$\blacktriangleright$}
\setbeamertemplate{itemize subitem}{\color{CalPolyGreen}$\bullet$}

\newcommand{\CourseNumber}{Course}
\newcommand{\CourseTitle}{Course Title}
\newcommand{\LectureNumber}{0}
\newcommand{\LectureTitle}{Lecture Title}
\newcommand{\LectureDate}{Date}
\newcommand{\CourseUnit}{Unit}

\newcommand{\coursenumber}[1]{\renewcommand{\CourseNumber}{#1}}
\newcommand{\coursetitle}[1]{\renewcommand{\CourseTitle}{#1}}
\newcommand{\lecturenumber}[1]{\renewcommand{\LectureNumber}{#1}}
\newcommand{\lecturetitle}[1]{\renewcommand{\LectureTitle}{#1}}
\newcommand{\lecturedate}[1]{\renewcommand{\LectureDate}{#1}}
\newcommand{\courseunit}[1]{\renewcommand{\CourseUnit}{#1}}

\author{Kirk Duran}
\institute{California Polytechnic State University, San Luis Obispo}
\date{\LectureDate}

\newcommand{\makedecktitle}{
  \begin{frame}[plain]
    \vspace{0.25in}
    {\usebeamerfont{title}\color{CalPolyGreen}\LectureTitle\par}
    \vspace{0.18in}
    {\large \CourseNumber{}: \CourseTitle\par}
    \vspace{0.08in}
    {\large Lecture \LectureNumber{} -- \LectureDate\par}
    \vfill
    {\small Kirk Duran \textbar{} Cal Poly, San Luis Obispo\par}
  \end{frame}
}

\newcommand{\placeholdernote}{
  \begin{block}{Placeholder}
    Replace these bullets with the final lecture material, examples, and in-class activities.
  \end{block}
}

\AtBeginSection[]{
  \begin{frame}{Roadmap}
    \tableofcontents[currentsection]
  \end{frame}
}


\usepackage{tikz}

\graphicspath{{../assets/lecture-09/}}

% If an image file is not yet available, skip the visual slot.
\newcommand{\lectureimage}[2][1.75in]{%
  \IfFileExists{../assets/lecture-09/#2}{%
    \includegraphics[width=\linewidth,height=#1,keepaspectratio]{#2}%
  }{}%
}

% Explicit placeholder for visuals/examples that will be added later.
\newcommand{\lectureplaceholder}[2][2.35in]{%
  \begin{center}
    \fbox{%
      \begin{minipage}[c][#1][c]{0.90\linewidth}
        \centering
        \small\textit{#2}
      \end{minipage}%
    }
  \end{center}%
}

\setbeamersize{text margin left=0.42in,text margin right=0.42in}

\coursenumber{CS 4880}
\coursetitle{Artificial Intelligence}
\lecturenumber{9}
\lecturetitle{Alpha--Beta Pruning and Depth-Limited Game Search}
\lecturedate{September 17, 2026}
\courseunit{Search and Problem Solving}

\begin{document}

% -----------------------------------------------------------------------------
% Source slide 1 -> New slide 1
% -----------------------------------------------------------------------------
\makedecktitle

% -----------------------------------------------------------------------------
% Source slide 2 -> New slide 2
% -----------------------------------------------------------------------------
\begin{frame}[shrink=18]{Deep Blue: Search Is Powerful, but Expensive}
  \begin{columns}[T,onlytextwidth]
    \begin{column}{0.55\textwidth}
      Deep Blue's success depended on searching enormous game trees, but exhaustive minimax is wasteful.
      \begin{itemize}
        \item Many legal moves are strategically irrelevant.
        \item Some branches can be proven unable to change the final decision.
        \item In chess, terminal states may be far beyond any practical search horizon.
      \end{itemize}
      \begin{block}{Two questions for this lecture}
        Can we \textbf{avoid searching irrelevant branches}, and can we \textbf{stop before terminal states} without giving up useful decisions?
      \end{block}
    \end{column}
    \begin{column}{0.39\textwidth}
      \lectureplaceholder[2.55in]{Visual placeholder: Deep Blue / chess search tree with a small highlighted region and many faded branches.}
    \end{column}
  \end{columns}
  \begin{keyidea}
    Practical game playing combines exact pruning with controlled approximations of search depth.
  \end{keyidea}
\end{frame}

\section{From Minimax to Pruning}

% -----------------------------------------------------------------------------
% Source slide 3 -> New slide 3
% -----------------------------------------------------------------------------
\begin{frame}[shrink=18]{Minimax Review: Start at the Leaves}
  \begin{columns}[T,onlytextwidth]
    \begin{column}{0.50\textwidth}
      Recall a depth-3 minimax tree with alternating MAX and MIN layers.
      \begin{itemize}
        \item Terminal utilities are measured from MAX's perspective.
        \item MAX nodes choose the largest child value.
        \item MIN nodes choose the smallest child value.
        \item Values are propagated from leaves back toward the root.
      \end{itemize}
      \begin{block}{Example tree}
        Leaf utilities, left to right: $-1,3,5,1,-6,-4,0,9$.
      \end{block}
    \end{column}
    \begin{column}{0.44\textwidth}
      \lectureplaceholder[2.70in]{Example placeholder: full binary minimax tree before any leaf has been evaluated. Highlight the first leaf, $-1$.}
    \end{column}
  \end{columns}
\end{frame}

% -----------------------------------------------------------------------------
% Source slide 4 -> New slide 4
% -----------------------------------------------------------------------------
\begin{frame}[shrink=18]{Minimax Review: First Leaf}
  \begin{columns}[T,onlytextwidth]
    \begin{column}{0.48\textwidth}
      The leftmost MAX node begins by evaluating its first child.
      \begin{itemize}
        \item First observed utility: $-1$.
        \item MAX's current best value is therefore $-1$.
        \item The sibling must still be inspected because it may be larger.
      \end{itemize}
      \begin{keyidea}
        A partial search gives a bound on what a node can currently guarantee, but not necessarily its final minimax value.
      \end{keyidea}
    \end{column}
    \begin{column}{0.46\textwidth}
      \lectureplaceholder[2.65in]{Example placeholder: same tree with leaf $-1$ revealed and the leftmost MAX node temporarily valued at $-1$.}
    \end{column}
  \end{columns}
\end{frame}

% -----------------------------------------------------------------------------
% Source slide 5 -> New slide 5
% -----------------------------------------------------------------------------
\begin{frame}[shrink=18]{Minimax Review: Finish the First MAX Node}
  \begin{columns}[T,onlytextwidth]
    \begin{column}{0.48\textwidth}
      The sibling leaf has utility $3$.
      \[
        \max(-1,3)=3.
      \]
      \begin{itemize}
        \item The leftmost MAX node returns $3$ to its parent MIN node.
        \item This establishes one candidate outcome for that MIN node.
      \end{itemize}
      \begin{block}{What MIN knows}
        MIN can already force an outcome no greater than $3$ if it chooses this child.
      \end{block}
    \end{column}
    \begin{column}{0.46\textwidth}
      \lectureplaceholder[2.65in]{Example placeholder: leaves $-1$ and $3$ revealed; their MAX parent labeled $3$.}
    \end{column}
  \end{columns}
\end{frame}

% -----------------------------------------------------------------------------
% Source slide 6 -> New slide 6
% -----------------------------------------------------------------------------
\begin{frame}[shrink=18]{Minimax Review: Finish the Left MIN Subtree}
  \begin{columns}[T,onlytextwidth]
    \begin{column}{0.50\textwidth}
      The second MAX child evaluates to
      \[
        \max(5,1)=5.
      \]
      The parent MIN node therefore returns
      \[
        \min(3,5)=3.
      \]
      \begin{itemize}
        \item MAX would prefer the branch worth $5$.
        \item MIN rejects it because the branch worth $3$ is better for MIN.
      \end{itemize}
    \end{column}
    \begin{column}{0.44\textwidth}
      \lectureplaceholder[2.65in]{Example placeholder: complete left half of tree, with MAX values $3$ and $5$, then MIN value $3$.}
    \end{column}
  \end{columns}
  \begin{keyidea}
    Minimax values encode what rational players can force, not what either player merely hopes will happen.
  \end{keyidea}
\end{frame}

% -----------------------------------------------------------------------------
% Source slide 7 -> New slide 7
% -----------------------------------------------------------------------------
\begin{frame}[shrink=18]{Minimax Review: Full Backed-Up Value}
  \begin{columns}[T,onlytextwidth]
    \begin{column}{0.50\textwidth}
      Completing the right subtree gives
      \[
        \max(-6,-4)=-4,\qquad \max(0,9)=9,
      \]
      and therefore
      \[
        \min(-4,9)=-4.
      \]
      At the root,
      \[
        \max(3,-4)=3.
      \]
      \begin{block}{Minimax decision}
        MAX chooses the left action because it guarantees utility $3$ under optimal opposition.
      \end{block}
    \end{column}
    \begin{column}{0.44\textwidth}
      \lectureplaceholder[2.70in]{Example placeholder: complete tree with internal values $3,5,-4,9$, MIN values $3,-4$, and root value $3$.}
    \end{column}
  \end{columns}
\end{frame}

% -----------------------------------------------------------------------------
% Source slide 8 -> New slide 8
% -----------------------------------------------------------------------------
\begin{frame}[shrink=18]{Why Full Minimax Becomes Impractical}
  \begin{columns}[T,onlytextwidth]
    \begin{column}{0.54\textwidth}
      Even a few plies of chess create a rapidly expanding set of positions.
      \begin{itemize}
        \item Each legal move creates a new state.
        \item The opponent then has many replies.
        \item Alternating choices multiply the number of candidate continuations.
        \item Most of those continuations will never occur in actual play.
      \end{itemize}
      \begin{block}{Search-tree growth}
        With branching factor $b$ and depth $m$, exhaustive minimax examines $O(b^m)$ nodes.
      \end{block}
    \end{column}
    \begin{column}{0.40\textwidth}
      \lectureplaceholder[2.60in]{Visual placeholder: four successive chess positions showing one line of play while many alternatives branch away.}
    \end{column}
  \end{columns}
\end{frame}

% -----------------------------------------------------------------------------
% Source slide 9 -> New slide 9
% -----------------------------------------------------------------------------
\begin{frame}[shrink=18]{The Cost of Looking Only Four Rounds Ahead}
  \begin{columns}[T,onlytextwidth]
    \begin{column}{0.54\textwidth}
      The source example expands a chess tree only a few rounds deep and already produces an enormous number of nodes.
      \begin{itemize}
        \item Brute-force minimax treats every legal continuation as potentially relevant.
        \item Yet many branches are dominated by alternatives already discovered elsewhere.
        \item If a branch cannot change the root decision, evaluating it is wasted work.
      \end{itemize}
      \begin{block}{Optimization target}
        Preserve the exact minimax result while examining fewer nodes.
      \end{block}
    \end{column}
    \begin{column}{0.40\textwidth}
      \lectureplaceholder[2.70in]{Visual placeholder: wide chess minimax tree after four rounds; annotate the source example's roughly 206,000 explored nodes.}
    \end{column}
  \end{columns}
\end{frame}

% -----------------------------------------------------------------------------
% Source slide 10 -> New slide 10
% -----------------------------------------------------------------------------
\begin{frame}[shrink=18]{Pruning Intuition: A Branch Already Looks Too Good for MAX}
  \begin{columns}[T,onlytextwidth]
    \begin{column}{0.50\textwidth}
      Return to the same minimax tree.
      \begin{itemize}
        \item The first MAX child under the left MIN node has value $3$.
        \item In the sibling MAX node, the first examined leaf has value $5$.
        \item Therefore that MAX node will return at least $5$, regardless of its unseen sibling.
      \end{itemize}
      \begin{block}{Observation}
        MIN already has a child worth $3$. Why would MIN ever choose a sibling guaranteed to be at least $5$?
      \end{block}
    \end{column}
    \begin{column}{0.44\textwidth}
      \lectureplaceholder[2.70in]{Example placeholder: left MIN subtree with known child value $3$ and sibling MAX node currently known to be $\ge 5$; hide the remaining leaf.}
    \end{column}
  \end{columns}
\end{frame}

% -----------------------------------------------------------------------------
% Source slide 11 -> New slide 11
% -----------------------------------------------------------------------------
\begin{frame}[shrink=18]{Pruning Intuition: MIN Will Never Choose That Branch}
  \begin{columns}[T,onlytextwidth]
    \begin{column}{0.52\textwidth}
      MIN prefers smaller values.
      \begin{itemize}
        \item MIN can already choose the branch worth $3$.
        \item The competing MAX branch is guaranteed to be at least $5$.
        \item Its unseen child cannot make that branch better for MIN.
      \end{itemize}
      \begin{block}{Safe conclusion}
        The remaining child of that MAX node can be ignored without changing the minimax value of the parent.
      \end{block}
    \end{column}
    \begin{column}{0.42\textwidth}
      \lectureplaceholder[2.70in]{Example placeholder: same tree with the unknown sibling visibly crossed out or faded; left MIN value remains $3$.}
    \end{column}
  \end{columns}
\end{frame}

% -----------------------------------------------------------------------------
% Source slide 12 -> New slide 12
% -----------------------------------------------------------------------------
\begin{frame}[shrink=18]{Pruning Intuition: The Symmetric Case}
  \begin{columns}[T,onlytextwidth]
    \begin{column}{0.52\textwidth}
      The same reasoning appears on the right side from MAX's perspective.
      \begin{itemize}
        \item The root already has a left action worth $3$.
        \item Suppose the right MIN node discovers a child worth $-4$.
        \item That MIN node can force a value no greater than $-4$.
        \item MAX would never choose that right action over the known value $3$.
      \end{itemize}
    \end{column}
    \begin{column}{0.42\textwidth}
      \lectureplaceholder[2.70in]{Example placeholder: root known to have value at least $3$ from its left child; right MIN subtree known to be $\le -4$ with an unexplored sibling.}
    \end{column}
  \end{columns}
  \begin{keyidea}
    Once an ancestor already has a better alternative, some descendants become irrelevant to the final decision.
  \end{keyidea}
\end{frame}

% -----------------------------------------------------------------------------
% Source slide 13 -> New slide 13
% -----------------------------------------------------------------------------
\begin{frame}[shrink=18]{Pruning Preserves the Same Root Decision}
  \begin{columns}[T,onlytextwidth]
    \begin{column}{0.50\textwidth}
      We have now skipped two portions of the tree.
      \begin{itemize}
        \item The pruned leaves are not approximations.
        \item Their exact values are unnecessary because they cannot affect the backed-up root value.
        \item The root still evaluates to $3$.
      \end{itemize}
      \begin{block}{Goal achieved}
        Search fewer nodes while returning exactly the same move as full minimax.
      \end{block}
    \end{column}
    \begin{column}{0.44\textwidth}
      \lectureplaceholder[2.70in]{Example placeholder: completed root value $3$ with pruned leaves faded; emphasize identical decision to full minimax.}
    \end{column}
  \end{columns}
\end{frame}

% -----------------------------------------------------------------------------
% Source slide 14 -> New slide 14
% -----------------------------------------------------------------------------
\begin{frame}[shrink=18]{How Much Can Pruning Save?}
  \begin{columns}[T,onlytextwidth]
    \begin{column}{0.53\textwidth}
      In the source demonstration, alpha--beta search at depth $4$ explored only a tiny fraction of the positions examined by naive minimax.
      \begin{itemize}
        \item Minimax: roughly $206{,}000$ positions in the illustrated tree.
        \item Alpha--beta: roughly $580$ positions in the illustrated comparison.
        \item Reported search fraction: about $0.28\%$ of the brute-force total.
      \end{itemize}
      \begin{block}{Important qualification}
        The savings depend strongly on the order in which moves are examined.
      \end{block}
    \end{column}
    \begin{column}{0.41\textwidth}
      \lectureplaceholder[2.70in]{Visual placeholder: side-by-side search trees labeled ``Minimax'' and ``Alpha--Beta,'' matching the source comparison.}
    \end{column}
  \end{columns}
\end{frame}

\section{Alpha--Beta Pruning}

% -----------------------------------------------------------------------------
% Source slide 15 -> New slide 15
% -----------------------------------------------------------------------------
\begin{frame}[shrink=18]{Alpha--Beta Pruning}
  \begin{cpdefinition}
    \textbf{Alpha--beta pruning} computes the same minimax value as exhaustive minimax while avoiding branches that cannot affect the final decision.
  \end{cpdefinition}
  \begin{columns}[T,onlytextwidth]
    \begin{column}{0.54\textwidth}
      The search maintains two bounds along the current path:
      \begin{itemize}
        \item $\alpha$: the best value MAX can guarantee so far.
        \item $\beta$: the best value MIN can guarantee so far.
      \end{itemize}
      \begin{block}{Cutoff condition}
        When $\alpha \ge \beta$, the remaining children of the current node cannot influence the final minimax choice and may be pruned.
      \end{block}
    \end{column}
    \begin{column}{0.40\textwidth}
      \lectureplaceholder[2.55in]{Visual placeholder: game tree with $\alpha$ and $\beta$ bounds along the path and scissors cutting an irrelevant subtree.}
    \end{column}
  \end{columns}
  \begin{keyidea}
    Alpha and beta summarize what the ancestors already know well enough to reject some unexplored alternatives.
  \end{keyidea}
\end{frame}

% -----------------------------------------------------------------------------
% Source slide 16 -> New slide 16
% -----------------------------------------------------------------------------
\begin{frame}[shrink=18]{Alpha--Beta Search Algorithm}
  \begin{columns}[T,onlytextwidth]
    \begin{column}{0.60\textwidth}
      \begin{block}{Recursive structure}
      \ttfamily\scriptsize
      ALPHA\_BETA(state, $\alpha$, $\beta$, maximizing)\\[0.02in]
      \quad if TERMINAL(state): return UTILITY(state)\\[0.04in]
      \quad if maximizing:\\
      \qquad value $\leftarrow -\infty$\\
      \qquad for action in ACTIONS(state):\\
      \qquad\quad value $\leftarrow \max$(value, ALPHA\_BETA(RESULT(...), $\alpha$, $\beta$, false))\\
      \qquad\quad $\alpha \leftarrow \max(\alpha,$ value$)$\\
      \qquad\quad if $\alpha \ge \beta$: break\\
      \qquad return value\\[0.04in]
      \quad else:\\
      \qquad value $\leftarrow +\infty$\\
      \qquad for action in ACTIONS(state):\\
      \qquad\quad value $\leftarrow \min$(value, ALPHA\_BETA(RESULT(...), $\alpha$, $\beta$, true))\\
      \qquad\quad $\beta \leftarrow \min(\beta,$ value$)$\\
      \qquad\quad if $\alpha \ge \beta$: break\\
      \qquad return value
      \end{block}
    \end{column}
    \begin{column}{0.34\textwidth}
      \begin{itemize}
        \item Same recursive minimax structure.
        \item Pass $\alpha$ and $\beta$ downward.
        \item Update the relevant bound after each child.
        \item Stop when the search window closes.
      \end{itemize}
      \lectureplaceholder[1.20in]{Code screenshot placeholder: Python alpha--beta implementation.}
    \end{column}
  \end{columns}
\end{frame}

\section{Worked Alpha--Beta Trace}

% -----------------------------------------------------------------------------
% Source slide 17 -> New slide 17
% -----------------------------------------------------------------------------
\begin{frame}[shrink=18]{Alpha--Beta Example: Initialize the Search}
  \begin{columns}[T,onlytextwidth]
    \begin{column}{0.48\textwidth}
      Start at the root MAX node with
      \[
        \alpha=-\infty,\qquad \beta=+\infty.
      \]
      These bounds are inherited by descendants until tighter information is discovered.
      \begin{block}{First step}
        Descend depth-first to the leftmost terminal leaf, whose utility is $-1$.
      \end{block}
    \end{column}
    \begin{column}{0.46\textwidth}
      \lectureplaceholder[2.70in]{Trace placeholder: initial alpha--beta tree; show $\alpha=-\infty,\beta=+\infty$ at root and highlight leaf $-1$.}
    \end{column}
  \end{columns}
\end{frame}

% -----------------------------------------------------------------------------
% Source slide 18 -> New slide 18
% -----------------------------------------------------------------------------
\begin{frame}[shrink=18]{Alpha--Beta Example: Update $\alpha$ at MAX}
  \begin{columns}[T,onlytextwidth]
    \begin{column}{0.48\textwidth}
      At the leftmost MAX node, the first child returns $-1$.
      \[
        \alpha \leftarrow \max(-\infty,-1)=-1.
      \]
      \begin{itemize}
        \item MAX can now guarantee at least $-1$ at this node.
        \item The sibling is still relevant because it may improve MAX's value.
      \end{itemize}
    \end{column}
    \begin{column}{0.46\textwidth}
      \lectureplaceholder[2.70in]{Trace placeholder: leftmost MAX node annotated $\alpha=-1,\beta=+\infty$ after leaf $-1$.}
    \end{column}
  \end{columns}
\end{frame}

% -----------------------------------------------------------------------------
% Source slide 19 -> New slide 19
% -----------------------------------------------------------------------------
\begin{frame}[shrink=18]{Alpha--Beta Example: Finish the First MAX Node}
  \begin{columns}[T,onlytextwidth]
    \begin{column}{0.48\textwidth}
      The sibling leaf returns $3$.
      \[
        \alpha \leftarrow \max(-1,3)=3.
      \]
      The MAX node is now complete and returns $3$ to its parent MIN node.
      \begin{block}{No cutoff yet}
        $\alpha=3$ and $\beta=+\infty$, so the current window remains open.
      \end{block}
    \end{column}
    \begin{column}{0.46\textwidth}
      \lectureplaceholder[2.70in]{Trace placeholder: leaves $-1,3$ evaluated; first MAX node labeled value $3$ and $\alpha=3$.}
    \end{column}
  \end{columns}
\end{frame}

% -----------------------------------------------------------------------------
% Source slide 20 -> New slide 20
% -----------------------------------------------------------------------------
\begin{frame}[shrink=18]{Alpha--Beta Example: Update $\beta$ at MIN}
  \begin{columns}[T,onlytextwidth]
    \begin{column}{0.49\textwidth}
      The parent is a MIN node. Its first child returns $3$.
      \[
        \beta \leftarrow \min(+\infty,3)=3.
      \]
      \begin{itemize}
        \item MIN can now guarantee a value no greater than $3$.
        \item This $\beta=3$ bound is passed to the next MAX child.
      \end{itemize}
    \end{column}
    \begin{column}{0.45\textwidth}
      \lectureplaceholder[2.70in]{Trace placeholder: left MIN node labeled $\alpha=-\infty,\beta=3$; move to its second MAX child.}
    \end{column}
  \end{columns}
\end{frame}

% -----------------------------------------------------------------------------
% Source slide 21 -> New slide 21
% -----------------------------------------------------------------------------
\begin{frame}[shrink=18]{Alpha--Beta Example: A Candidate Already Exceeds $\beta$}
  \begin{columns}[T,onlytextwidth]
    \begin{column}{0.49\textwidth}
      The second MAX child is searched under the window
      \[
        \alpha=-\infty,\qquad \beta=3.
      \]
      Its first leaf has value $5$, so
      \[
        \alpha \leftarrow 5.
      \]
      \begin{block}{Critical comparison}
        Now $\alpha=5 \ge \beta=3$.
      \end{block}
    \end{column}
    \begin{column}{0.45\textwidth}
      \lectureplaceholder[2.70in]{Trace placeholder: second MAX node with first leaf $5$ revealed; annotate $\alpha=5,\beta=3$.}
    \end{column}
  \end{columns}
\end{frame}

% -----------------------------------------------------------------------------
% Source slide 22 -> New slide 22
% -----------------------------------------------------------------------------
\begin{frame}[shrink=18]{Alpha--Beta Example: First Cutoff}
  \begin{columns}[T,onlytextwidth]
    \begin{column}{0.50\textwidth}
      Because $\alpha \ge \beta$, the remaining child is irrelevant.
      \begin{itemize}
        \item This MAX node is already worth at least $5$.
        \item Its parent MIN node already has an alternative worth $3$.
        \item MIN will never select a branch that cannot fall below $5$.
      \end{itemize}
      \begin{block}{Prune}
        Do not evaluate the sibling leaf $1$.
      \end{block}
    \end{column}
    \begin{column}{0.44\textwidth}
      \lectureplaceholder[2.70in]{Trace placeholder: show cutoff at $\alpha=5,\beta=3$ and cross out the unseen leaf $1$.}
    \end{column}
  \end{columns}
  \begin{keyidea}
    The cutoff is justified by ancestor preferences, not by guessing the value of the pruned node.
  \end{keyidea}
\end{frame}

% -----------------------------------------------------------------------------
% Source slide 23 -> New slide 23
% -----------------------------------------------------------------------------
\begin{frame}[shrink=18]{Alpha--Beta Example: Return the Left Subtree Value}
  \begin{columns}[T,onlytextwidth]
    \begin{column}{0.50\textwidth}
      The left MIN node compares its children:
      \begin{itemize}
        \item first child: exactly $3$;
        \item second child: known to be at least $5$.
      \end{itemize}
      Therefore MIN returns $3$.
      \[
        \min(3,\ge 5)=3.
      \]
      \begin{block}{At the root}
        MAX updates its bound to $\alpha=3$.
      \end{block}
    \end{column}
    \begin{column}{0.44\textwidth}
      \lectureplaceholder[2.70in]{Trace placeholder: left MIN subtree returns $3$; pruned leaf remains faded. Root now records $\alpha=3$.}
    \end{column}
  \end{columns}
\end{frame}

% -----------------------------------------------------------------------------
% Source slide 24 -> New slide 24
% -----------------------------------------------------------------------------
\begin{frame}[shrink=18]{Alpha--Beta Example: Search the Right MIN Subtree}
  \begin{columns}[T,onlytextwidth]
    \begin{column}{0.50\textwidth}
      The root enters its right MIN child with
      \[
        \alpha=3,\qquad \beta=+\infty.
      \]
      The first MAX grandchild evaluates leaves $-6$ and $-4$:
      \[
        \max(-6,-4)=-4.
      \]
      This value is returned to the right MIN node.
    \end{column}
    \begin{column}{0.44\textwidth}
      \lectureplaceholder[2.70in]{Trace placeholder: right MIN subtree active; first MAX child evaluates $-6,-4$ and returns $-4$ while root bound $\alpha=3$ is visible.}
    \end{column}
  \end{columns}
\end{frame}

% -----------------------------------------------------------------------------
% Source slide 25 -> New slide 25
% -----------------------------------------------------------------------------
\begin{frame}[shrink=18]{Alpha--Beta Example: Update $\beta$ to $-4$}
  \begin{columns}[T,onlytextwidth]
    \begin{column}{0.50\textwidth}
      At the right MIN node,
      \[
        \beta \leftarrow \min(+\infty,-4)=-4.
      \]
      But this subtree was entered with $\alpha=3$ from the root.
      \begin{block}{Window closes}
        $\alpha=3 \ge \beta=-4$.
      \end{block}
      MIN can already force $-4$, which is worse for MAX than the root's known alternative $3$.
    \end{column}
    \begin{column}{0.44\textwidth}
      \lectureplaceholder[2.70in]{Trace placeholder: right MIN node labeled $\alpha=3,\beta=-4$ after its first child returns $-4$.}
    \end{column}
  \end{columns}
\end{frame}

% -----------------------------------------------------------------------------
% Source slide 26 -> New slide 26
% -----------------------------------------------------------------------------
\begin{frame}[shrink=18]{Alpha--Beta Example: Second Cutoff and Final Decision}
  \begin{columns}[T,onlytextwidth]
    \begin{column}{0.50\textwidth}
      The rest of the right MIN subtree can be pruned.
      \begin{itemize}
        \item The unseen MAX child containing $0$ and $9$ is never evaluated.
        \item The right root action is known to be no better than $-4$.
        \item The root already has an action worth $3$.
      \end{itemize}
      \begin{block}{Final result}
        Alpha--beta returns root value $3$, exactly matching full minimax.
      \end{block}
    \end{column}
    \begin{column}{0.44\textwidth}
      \lectureplaceholder[2.70in]{Trace placeholder: final alpha--beta tree with both cutoffs visible and root value $3$.}
    \end{column}
  \end{columns}
\end{frame}

% -----------------------------------------------------------------------------
% Source slide 27 -> New slide 27
% -----------------------------------------------------------------------------
\begin{frame}[shrink=18]{The Pruning Logic}
  \begin{columns}[T,onlytextwidth]
    \begin{column}{0.47\textwidth}
      \begin{block}{At a MAX node}
        \begin{itemize}
          \item MAX searches for the highest value.
          \item $\alpha$ tracks MAX's best guarantee so far.
          \item If $\alpha \ge \beta$, an ancestor MIN node already has a better alternative.
        \end{itemize}
      \end{block}
    \end{column}
    \begin{column}{0.47\textwidth}
      \begin{block}{At a MIN node}
        \begin{itemize}
          \item MIN searches for the lowest value.
          \item $\beta$ tracks MIN's best guarantee so far.
          \item If $\beta \le \alpha$, an ancestor MAX node already has a better alternative.
        \end{itemize}
      \end{block}
    \end{column}
  \end{columns}
  \lectureplaceholder[1.25in]{Visual placeholder: compact annotated tree showing one MAX cutoff and one MIN cutoff.}
  \begin{keyidea}
    A cutoff occurs when the current branch falls outside the range of values that could still matter to an ancestor.
  \end{keyidea}
\end{frame}

% -----------------------------------------------------------------------------
% Source slide 28 -> New slide 28
% -----------------------------------------------------------------------------
\begin{frame}[shrink=18]{Alpha--Beta Properties}
  \begin{center}
    \renewcommand{\arraystretch}{1.35}
    \begin{tabular}{l|p{0.66\textwidth}}
      \textbf{Property} & \textbf{Alpha--beta search} \\\hline
      Complete & Yes, for finite game trees. \\
      Optimal & Yes, under the same optimal-play assumptions as minimax. \\
      Worst-case time & $O(b^m)$, the same asymptotic bound as minimax. \\
      Best-case time & $O(b^{m/2})$ with near-perfect move ordering. \\
      Space & $O(m)$ for depth-first recursive generation, ignoring stored move lists. \\
    \end{tabular}
  \end{center}
  \begin{keyidea}
    Alpha--beta improves efficiency without changing the minimax decision; move ordering determines how much pruning actually occurs.
  \end{keyidea}
\end{frame}

\section{Depth Limits and Evaluation Functions}

% -----------------------------------------------------------------------------
% Source slide 29 -> New slide 29
% -----------------------------------------------------------------------------
\begin{frame}[shrink=18]{Pruning Is Not Enough}
  \begin{columns}[T,onlytextwidth]
    \begin{column}{0.54\textwidth}
      Even after alpha--beta pruning, chess trees remain extremely deep.
      \begin{itemize}
        \item A game may require dozens of moves before reaching checkmate, resignation, or draw.
        \item Some variations are much deeper than others.
        \item Waiting for terminal utilities is therefore unrealistic under fixed time limits.
      \end{itemize}
      \begin{block}{New problem}
        If search stops before the game ends, what value should be assigned to a nonterminal position?
      \end{block}
    \end{column}
    \begin{column}{0.40\textwidth}
      \lectureplaceholder[2.60in]{Visual placeholder: Deep Blue / chessboard with an uneven deep tree extending far below the current search frontier.}
    \end{column}
  \end{columns}
\end{frame}

% -----------------------------------------------------------------------------
% Source slide 30 -> New slide 30
% -----------------------------------------------------------------------------
\begin{frame}[shrink=18]{Depth-Limited Minimax}
  \begin{cpdefinition}
    \textbf{Depth-limited minimax} stops search at a cutoff depth and evaluates the resulting nonterminal states with an evaluation function.
  \end{cpdefinition}
  \begin{columns}[T,onlytextwidth]
    \begin{column}{0.55\textwidth}
      Replace the terminal-only base case with two possibilities:
      \begin{itemize}
        \item if the state is terminal, return its true utility;
        \item if the cutoff condition is reached, return $\mathrm{EVAL}(s)$.
      \end{itemize}
      The returned estimates are then backed up through MAX and MIN exactly as before.
    \end{column}
    \begin{column}{0.39\textwidth}
      \lectureplaceholder[2.45in]{Visual placeholder: game tree truncated by a horizontal ``search horizon''; frontier nodes labeled with heuristic values.}
    \end{column}
  \end{columns}
  \begin{keyidea}
    Depth limiting trades exact terminal reasoning for a heuristic estimate of positions near the search horizon.
  \end{keyidea}
\end{frame}

% -----------------------------------------------------------------------------
% Source slide 31 -> New slide 31
% -----------------------------------------------------------------------------
\begin{frame}[shrink=18]{Depth Limits and Pruning Work Together}
  \begin{columns}[T,onlytextwidth]
    \begin{column}{0.53\textwidth}
      Practical game search usually combines both mechanisms.
      \begin{itemize}
        \item \textbf{Depth limit}: controls how far the tree may grow.
        \item \textbf{Alpha--beta}: avoids searching branches within that horizon that cannot matter.
        \item \textbf{Evaluation function}: estimates frontier states when terminal utilities are unavailable.
      \end{itemize}
      \begin{block}{Source comparison}
        The lecture's depth-4 example reports alpha--beta examining about $0.28\%$ of the nodes shown for brute-force minimax.
      \end{block}
    \end{column}
    \begin{column}{0.41\textwidth}
      \lectureplaceholder[2.70in]{Visual placeholder: repeat source side-by-side minimax vs. alpha--beta depth-4 trees, now annotate the shared cutoff depth.}
    \end{column}
  \end{columns}
\end{frame}

% -----------------------------------------------------------------------------
% Source slide 32 -> New slide 32
% -----------------------------------------------------------------------------
\begin{frame}[shrink=18]{Depth-Limited Alpha--Beta Search}
  \begin{columns}[T,onlytextwidth]
    \begin{column}{0.62\textwidth}
      \begin{block}{Algorithmic structure}
      \ttfamily\scriptsize
      ALPHA\_BETA(state, depth, $\alpha$, $\beta$, maximizing)\\[0.02in]
      \quad if TERMINAL(state): return UTILITY(state)\\
      \quad if CUTOFF(state, depth): return EVAL(state)\\[0.04in]
      \quad if maximizing:\\
      \qquad value $\leftarrow -\infty$\\
      \qquad for action in ACTIONS(state):\\
      \qquad\quad value $\leftarrow \max$(value, ALPHA\_BETA(RESULT(...), depth$+1$, $\alpha$, $\beta$, false))\\
      \qquad\quad $\alpha \leftarrow \max(\alpha,$ value$)$\\
      \qquad\quad if $\alpha \ge \beta$: break\\
      \qquad return value\\[0.04in]
      \quad else: \emph{symmetrically minimize and update $\beta$}
      \end{block}
    \end{column}
    \begin{column}{0.32\textwidth}
      \begin{itemize}
        \item Track current depth.
        \item Stop at terminal or cutoff states.
        \item Evaluate frontier states heuristically.
        \item Preserve alpha--beta pruning above the frontier.
      \end{itemize}
      \lectureplaceholder[1.05in]{Code screenshot placeholder: depth-limited alpha--beta.}
    \end{column}
  \end{columns}
\end{frame}

% -----------------------------------------------------------------------------
% Source slide 33 -> New slide 33
% -----------------------------------------------------------------------------
\begin{frame}[shrink=18]{Evaluation Functions Replace Unreachable Utilities}
  \begin{columns}[T,onlytextwidth]
    \begin{column}{0.55\textwidth}
      When search stops before terminal states, the agent needs a proxy for true utility.
      \begin{itemize}
        \item The evaluation should correlate with the probability or quality of eventual outcomes.
        \item Positive values should favor MAX; negative values should favor MIN.
        \item Evaluation must be cheap enough to apply at many frontier nodes.
      \end{itemize}
      \begin{block}{Design objective}
        Estimate the strategic value of a state using features available now, without solving the remainder of the game.
      \end{block}
    \end{column}
    \begin{column}{0.39\textwidth}
      \lectureplaceholder[2.55in]{Visual placeholder: Deep Blue examining a frontier of chess states, each assigned a heuristic score.}
    \end{column}
  \end{columns}
\end{frame}

% -----------------------------------------------------------------------------
% Source slide 34 -> New slide 34
% -----------------------------------------------------------------------------
\begin{frame}[shrink=18]{Chess Evaluation Functions}
  \begin{columns}[T,onlytextwidth]
    \begin{column}{0.55\textwidth}
      A classical evaluation function combines weighted strategic features:
      \[
        \mathrm{EVAL}(s)=w_1f_1(s)+w_2f_2(s)+\cdots+w_nf_n(s).
      \]
      Typical features include
      \begin{itemize}
        \item material balance;
        \item king safety;
        \item pawn structure;
        \item mobility and control of important squares;
        \item piece activity and positional constraints.
      \end{itemize}
    \end{column}
    \begin{column}{0.39\textwidth}
      \lectureplaceholder[2.55in]{Example placeholder: chess position annotated with material, king-safety, pawn-structure, and mobility features.}
    \end{column}
  \end{columns}
  \begin{keyidea}
    An evaluation function is useful only insofar as its features predict long-term game outcomes beyond the current search depth.
  \end{keyidea}
\end{frame}

% -----------------------------------------------------------------------------
% Source slide 35 -> New slide 35
% -----------------------------------------------------------------------------
\begin{frame}[shrink=18]{Connect Four: A Simple Heuristic Evaluation}
  \begin{columns}[T,onlytextwidth]
    \begin{column}{0.54\textwidth}
      A Connect Four heuristic can score patterns that indicate progress toward a win.
      \begin{itemize}
        \item Large positive score for four in a row.
        \item Smaller positive score for three in a row with an open cell.
        \item Smaller still for two in a row with multiple open cells.
        \item Apply corresponding negative values to dangerous opponent patterns.
      \end{itemize}
      \begin{block}{Example scoring sketch}
        $+1000$ for four; $+100$ for an open three; $+10$ for an open two, with analogous penalties for the opponent.
      \end{block}
    \end{column}
    \begin{column}{0.40\textwidth}
      \lectureplaceholder[2.55in]{Example placeholder: Connect Four board with highlighted windows/patterns and a small scoring-function code excerpt.}
    \end{column}
  \end{columns}
\end{frame}

\section{Depth-Limit Failure Modes and Extensions}

% -----------------------------------------------------------------------------
% Source slide 36 -> New slide 36
% -----------------------------------------------------------------------------
\begin{frame}[shrink=18]{The Horizon Effect}
  \begin{cpdefinition}
    The \textbf{horizon effect} occurs when important consequences lie just beyond the search cutoff, causing the evaluation function to misjudge the current position.
  \end{cpdefinition}
  \begin{columns}[T,onlytextwidth]
    \begin{column}{0.54\textwidth}
      \begin{itemize}
        \item A tactic may look favorable only because its refutation is one ply beyond the horizon.
        \item An agent may delay an unavoidable loss until after the cutoff.
        \item The frontier position appears stable even though a forcing sequence is imminent.
      \end{itemize}
      \begin{block}{Consequence}
        A deeper fixed cutoff helps, but does not eliminate the problem; the tactical event can simply move beyond the new horizon.
      \end{block}
    \end{column}
    \begin{column}{0.40\textwidth}
      \lectureplaceholder[2.10in]{Visual placeholder: horizon line with a move that appears strong before the line but loses immediately beyond it.}
    \end{column}
  \end{columns}
\end{frame}

% -----------------------------------------------------------------------------
% Source slide 37 -> New slide 37
% -----------------------------------------------------------------------------
\begin{frame}[shrink=18]{Stable vs. Unstable Frontier Positions}
  \begin{columns}[T,onlytextwidth]
    \begin{column}{0.47\textwidth}
      \begin{block}{Quiescent position}
        A strategically stable position with no immediate forcing tactical sequence.
        \begin{itemize}
          \item no urgent captures;
          \item no immediate checks;
          \item no major promotions or exchanges pending.
        \end{itemize}
      \end{block}
    \end{column}
    \begin{column}{0.47\textwidth}
      \begin{block}{Non-quiescent position}
        An unstable position in which the evaluation may change sharply after a few forcing moves.
        \begin{itemize}
          \item tactical captures;
          \item checks or threats;
          \item promotions;
          \item unresolved material swings.
        \end{itemize}
      \end{block}
    \end{column}
  \end{columns}
  \begin{keyidea}
    Evaluation is most trustworthy when the cutoff occurs in a relatively stable position.
  \end{keyidea}
\end{frame}

% -----------------------------------------------------------------------------
% Source slide 38 -> New slide 38
% -----------------------------------------------------------------------------
\begin{frame}[shrink=18]{Quiescence Search}
  \begin{cpdefinition}
    \textbf{Quiescence search} selectively extends search beyond the nominal depth limit when the frontier position is tactically unstable.
  \end{cpdefinition}
  \begin{columns}[T,onlytextwidth]
    \begin{column}{0.54\textwidth}
      Rather than considering every legal move, continue only through forcing moves such as
      \begin{itemize}
        \item captures;
        \item checks;
        \item promotions;
        \item other domain-specific tactical actions.
      \end{itemize}
      Stop once a sufficiently stable state is reached and then apply the evaluation function.
    \end{column}
    \begin{column}{0.40\textwidth}
      \lectureplaceholder[2.50in]{Visual placeholder: main alpha--beta tree ending at a horizon, with a few selectively extended tactical branches below it.}
    \end{column}
  \end{columns}
  \begin{keyidea}
    Quiescence search spends extra depth where the evaluation function is least reliable.
  \end{keyidea}
\end{frame}

% -----------------------------------------------------------------------------
% Source slide 39 -> New slide 39
% -----------------------------------------------------------------------------
\begin{frame}[shrink=18]{Singular Extensions}
  \begin{columns}[T,onlytextwidth]
    \begin{column}{0.54\textwidth}
      A \textbf{singular extension} searches deeper along a line when one move appears substantially stronger than all alternatives.
      \begin{itemize}
        \item Trigger when the best move is clearly separated from the rest.
        \item Allocate additional depth to that critical continuation.
        \item Avoid spending equal depth on obviously inferior alternatives.
      \end{itemize}
      \begin{block}{Why it helps}
        Some strategically decisive sequences are narrow rather than broad; selective depth can reveal their true consequences.
      \end{block}
    \end{column}
    \begin{column}{0.40\textwidth}
      \lectureplaceholder[2.55in]{Visual placeholder: tree with one brightly highlighted branch extended several plies deeper than its siblings.}
    \end{column}
  \end{columns}
\end{frame}

% -----------------------------------------------------------------------------
% Source slide 40 -> New slide 40
% -----------------------------------------------------------------------------
\begin{frame}[shrink=18]{Forward Pruning: Search Less by Taking Risk}
  \begin{columns}[T,onlytextwidth]
    \begin{column}{0.48\textwidth}
      \textbf{Forward pruning} deliberately discards or reduces branches before their irrelevance is proven.
      \begin{itemize}
        \item Reduces search effort aggressively.
        \item Can reach greater effective depth.
        \item Sacrifices the minimax guarantee if a discarded move was actually important.
      \end{itemize}
    \end{column}
    \begin{column}{0.46\textwidth}
      \begin{block}{Examples from the source}
        \begin{itemize}
          \item Beam-style move selection.
          \item ProbCut / probability-based pruning.
          \item Late-move reductions (LMR).
        \end{itemize}
      \end{block}
      \lectureplaceholder[1.20in]{Visual placeholder: only a few promising branches remain bright while others are intentionally discarded.}
    \end{column}
  \end{columns}
  \begin{keyidea}
    Alpha--beta pruning is exact; forward-pruning methods are selective approximations that trade certainty for speed.
  \end{keyidea}
\end{frame}

% -----------------------------------------------------------------------------
% Source slide 41 -> New slide 41
% -----------------------------------------------------------------------------
\begin{frame}[shrink=18]{Search vs. Lookup}
  \begin{columns}[T,onlytextwidth]
    \begin{column}{0.47\textwidth}
      \begin{block}{Opening knowledge}
        Familiar early-game positions can be stored in opening books rather than rediscovered through search every game.
      \end{block}
      \begin{itemize}
        \item curated or statistically derived lines;
        \item immediate access to known strong moves;
        \item saves computation for unfamiliar positions.
      \end{itemize}
    \end{column}
    \begin{column}{0.47\textwidth}
      \begin{block}{Endgame knowledge}
        Small endgames can be solved exhaustively in advance and stored in tablebases.
      \end{block}
      \begin{itemize}
        \item exact win/draw/loss information;
        \item often exact distance-to-conversion information;
        \item replace online search with direct lookup when applicable.
      \end{itemize}
    \end{column}
  \end{columns}
  \lectureplaceholder[1.20in]{Visual placeholder: opening book on the left, endgame tablebase/database on the right, search in the middle.}
\end{frame}

% -----------------------------------------------------------------------------
% Source slide 42 -> New slide 42
% -----------------------------------------------------------------------------
\begin{frame}[shrink=18]{Iterative Deepening Under a Time Limit}
  \begin{cpdefinition}
    \textbf{Iterative deepening} repeatedly performs depth-limited search with increasing depth until the available time is exhausted.
  \end{cpdefinition}
  \begin{columns}[T,onlytextwidth]
    \begin{column}{0.53\textwidth}
      \begin{enumerate}
        \item Search to depth $1$.
        \item Repeat at depths $2,3,\ldots$.
        \item Stop when the time budget is nearly exhausted.
        \item Return the best move from the deepest fully completed iteration.
      \end{enumerate}
    \end{column}
    \begin{column}{0.41\textwidth}
      \begin{block}{Why repeat work?}
        Earlier iterations improve move ordering for deeper alpha--beta searches and guarantee an anytime answer if computation must stop unexpectedly.
      \end{block}
      \lectureplaceholder[1.05in]{Visual placeholder: nested depth-1, depth-2, depth-3 search horizons around a chess clock.}
    \end{column}
  \end{columns}
\end{frame}

\section{Putting the Mechanisms Together}

% -----------------------------------------------------------------------------
% Source slide 43 -> New slide 43
% -----------------------------------------------------------------------------
\begin{frame}[shrink=18]{Deep Blue as a Search System}
  \begin{columns}[T,onlytextwidth]
    \begin{column}{0.56\textwidth}
      The source lecture uses Deep Blue to connect the mechanisms into one practical architecture:
      \begin{itemize}
        \item minimax-style adversarial reasoning;
        \item alpha--beta pruning;
        \item depth-limited search;
        \item heuristic evaluation of nonterminal states;
        \item quiescence search and selective extensions;
        \item forward-pruning ideas;
        \item opening books and endgame knowledge;
        \item progressive / iterative deepening.
      \end{itemize}
    \end{column}
    \begin{column}{0.38\textwidth}
      \lectureplaceholder[2.60in]{Visual placeholder: Deep Blue match photo plus a layered pipeline from search to pruning to evaluation to lookup.}
    \end{column}
  \end{columns}
  \begin{keyidea}
    Strong game-playing systems are not one algorithm; they are carefully engineered combinations of exact search, heuristic evaluation, selective depth, and stored knowledge.
  \end{keyidea}
\end{frame}

% -----------------------------------------------------------------------------
% Source slide 44 -> New slide 44
% -----------------------------------------------------------------------------
\begin{frame}[shrink=18]{Two Modern Paths to Strong Chess AI}
  \begin{columns}[T,onlytextwidth]
    \begin{column}{0.47\textwidth}
      \begin{block}{Search-engine tradition}
        Systems such as Stockfish build on highly optimized tree search.
        \begin{itemize}
          \item alpha--beta-style pruning;
          \item transposition tables;
          \item selective reductions and extensions;
          \item strong evaluation functions;
          \item opening and endgame knowledge.
        \end{itemize}
      \end{block}
    \end{column}
    \begin{column}{0.47\textwidth}
      \begin{block}{Learning-centered tradition}
        Systems such as AlphaZero combine learned value/policy models with Monte Carlo Tree Search.
        \begin{itemize}
          \item self-play learning;
          \item learned state evaluation;
          \item policy-guided search;
          \item no classical alpha--beta requirement.
        \end{itemize}
      \end{block}
    \end{column}
  \end{columns}
  \begin{keyidea}
    The core problem remains selective lookahead: devote computation to the continuations most likely to matter.
  \end{keyidea}
\end{frame}

% -----------------------------------------------------------------------------
% Source slide 45 -> New slide 45
% -----------------------------------------------------------------------------
\begin{frame}[shrink=18]{Chess AI: Comparing Machine and Human Strength}
  \begin{columns}[T,onlytextwidth]
    \begin{column}{0.55\textwidth}
      The source slide compares top human ratings with contemporary engine ratings to emphasize the scale of modern chess-engine performance.
      \begin{itemize}
        \item Engine strength now substantially exceeds elite human play under standard engine-testing conditions.
        \item Exact numerical rankings change over time and depend on the rating pool and hardware configuration.
        \item For lecture use, the conceptual point is more durable than any one snapshot.
      \end{itemize}
      \begin{block}{Update placeholder}
        Replace with a current human-vs.-engine rating snapshot before presenting if numerical rankings matter.
      \end{block}
    \end{column}
    \begin{column}{0.39\textwidth}
      \lectureplaceholder[2.55in]{Visual placeholder: current top-human rating list beside a current computer-chess engine rating list.}
    \end{column}
  \end{columns}
\end{frame}

% -----------------------------------------------------------------------------
% Source slide 46 -> New slide 46
% -----------------------------------------------------------------------------
\begin{frame}[shrink=18]{Real Games Violate the Idealized Model}
  \begin{columns}[T,onlytextwidth]
    \begin{column}{0.47\textwidth}
      \begin{block}{Assumptions so far}
        \begin{itemize}
          \item opponent plays optimally;
          \item legal actions are known;
          \item state is fully observable;
          \item action outcomes are deterministic;
          \item search can model the relevant future.
        \end{itemize}
      \end{block}
    \end{column}
    \begin{column}{0.47\textwidth}
      \begin{block}{Complications in real games}
        \begin{itemize}
          \item opponents can be suboptimal or adaptive;
          \item chance may affect outcomes;
          \item information may be hidden;
          \item exact outcomes may be probabilistic;
          \item the full tree may be too large even for sophisticated pruning.
        \end{itemize}
      \end{block}
    \end{column}
  \end{columns}
  \begin{keyidea}
    Alpha--beta is extremely effective for deterministic adversarial trees, but other game models require different search machinery.
  \end{keyidea}
\end{frame}

% -----------------------------------------------------------------------------
% Source slide 47 -> New slide 47
% -----------------------------------------------------------------------------
\begin{frame}[shrink=18]{Takeaways and Next Time}
  \begin{itemize}
    \item Alpha--beta pruning returns the same move as minimax while skipping branches that cannot affect the decision.
    \item $\alpha$ records MAX's best guarantee; $\beta$ records MIN's best guarantee; prune when $\alpha\ge\beta$.
    \item Move ordering controls how much pruning occurs: the best case reduces effective search from $O(b^m)$ toward $O(b^{m/2})$.
    \item Depth-limited search replaces unreachable terminal utilities with heuristic evaluation at a cutoff frontier.
    \item Horizon effects motivate quiescence search and selective extensions near tactically unstable positions.
    \item Iterative deepening provides time-bounded, anytime search while improving move ordering for alpha--beta.
    \item Next lecture: \textbf{Monte Carlo Tree Search}.
  \end{itemize}
  \begin{keyidea}
    Practical adversarial search succeeds by deciding both \emph{which branches deserve exploration} and \emph{how deeply each branch should be explored}.
  \end{keyidea}
\end{frame}

\end{document}
